\begin{textsample}{POS Dim 1 – df – Score 32.00 – df/df\_em\_20.txt}  \label{ex:f1_pos_097}
5.2 MATEMÁTICA E SUAS TECNOLOGIAS - FORMAÇÃO GERAL BÁSICA A sociedade e suas constantes mudanças \textbf{influenciam} diversas áreas do conhecimento, e com a Matemática não seria diferente.
Ernest ( 1991, apud D’Ambrosio, 1993 ) destaca o papel importante \textbf{que} a interação social \textbf{desempenha} na origem da Matemática.
Ele ressalta \textbf{que} “ a Matemática evolui através de \textbf{um} processo humano de criativo de geração de ideias e subsequente processo social de negociação de significados, simbolização, refutação e formalização ”.
Além disso, é proposto pelo \textbf{autor} \textbf{que} o conhecimento matemático, em seu princípio, “ evolui da resolução de problemas provenientes da realidade ou da própria construção matemática ”.
Essa preocupação em relação à compreensão da Matemática e, principalmente, sobre o modo como é trabalhado o seu ensino, faz -se presente em nossa sociedade há algumas décadas.
Ao analisar a trajetória da educação Matemática, é possível perceber \textbf{que} em certas décadas, como as de sessenta e setenta, a característica \textbf{predominante} era a de \textbf{um} ensino mais transmissivo, \textbf{que} praticamente desconsiderava as ações dos estudantes perante os conteúdos apresentados e se pautava no formalismo dos algoritmos e das fórmulas ( TAROUCO, SILVA, DA SILVA, 2016 ).
Desde essa época, o ensino de Matemática vem sendo questionado e, cada vez mais, busca -se aprimorá -lo.
Em relação a isso, os Parâmetros Curriculares Nacionais para o Ensino Médio ( PCNEM ) alertam para a necessidade de desenvolvimento de instrumentos \textbf{que} possibilitem ao estudante construir “ efetivamente as abstrações matemáticas, evitando -se a \textbf{memorização} indiscriminada de algoritmos, de forma prejudicial ao aprendizado ” ( BRASIL, 2000 ).
Em busca desse aprimoramento, não só do ensino da Matemática, mas do Ensino Médio de modo geral, os documentos oficiais, desde a Lei de Diretrizes e Bases da Educação Nacional ( LDB ), de 1996, até a Base Nacional Comum Curricular ( BNCC ), de 2018, apresentam premissas para reformular essa etapa com o objetivo de atender, cada vez mais, à necessidade de atualização da educação brasileira.
No \textbf{que} se refere à atualização da educação, D’Ambrosio faz \textbf{um} comparativo entre a Arte e a Matemática. > O importante relacionamento de Arte e Matemática não se faz com as músicas \textbf{que} eram \textbf{populares} no \textbf{século} XIX, valsas e minuetos, nem com pintura, escultura, arquitetura e literatura dessa época.
A arte do \textbf{século} XXI para os jovens do \textbf{século} XXI está intimamente relacionada com a matemática do \textbf{século} XXI, \textbf{que} é muito diferente da matemática do \textbf{século} XIX ( D’AMBROSIO, 2001 ).
Nesse trecho o \textbf{autor} \textbf{lembra} da importância de utilizar metodologias \textbf{atuais} para ensinar Matemática.
Em relação a isso, as Orientações Educacionais Complementares aos Parâmetros Curriculares Nacionais ( PCNEM+ ) da área de Ciências da Natureza, Matemática e Suas Tecnologias afirmam \textbf{que} o Ensino Médio se modifica de \textbf{um} \textbf{simples} preparatório para o ensino superior ou estritamente profissionalizante, para assumir a característica de completar a Educação Básica no \textbf{que} \textbf{diz} respeito a preparar o estudante “ para a vida, \textbf{qualificar} para a cidadania e capacitar para o aprendizado permanente, em eventual prosseguimento dos estudos ou diretamente no mundo do trabalho ” ( BRASIL, 2002a, p. 8 ).
Diante desse cenário, o Currículo em Movimento propõe, para a área da Matemática, a mesma linha de formação apresentada pela BNCC, ou seja, a construção de \textbf{uma} visão integrada da Matemática, aplicada à realidade, em diferentes contextos ( BRASIL, 2018a, p. 471 ).
Sob essa \textbf{ótica}, o estudante passa a ser vislumbrado como sujeito de sua prática cognitiva, dotado de conhecimento de mundo e de bagagem cultural, aspectos \textbf{que} impactam sobremaneira o processo de aprendizagem.
Portanto, tais conhecimentos devem ser aproveitados e entrelaçados aos objetivos de aprendizagem curriculares.
Nesse contexto, as Diretrizes Curriculares Nacionais para o Ensino Médio estabelecem \textbf{que} : > [... ] o currículo é conceituado como a proposta de ação educativa constituída pela seleção de conhecimentos construídos pela sociedade, expressando -se por práticas escolares \textbf{que} se desdobram em torno de conhecimentos relevantes e pertinentes, \textbf{permeadas} pelas relações sociais, articulando vivências e saberes dos estudantes e contribuindo para o desenvolvimento de suas identidades e condições cognitivas e socioemocionais ( BRASIL, 2018b ).
Sendo assim, este Currículo assume a responsabilidade de firmar todo o potencial já desenvolvido pelos estudantes no Ensino Fundamental, com o intuito de impulsionar ações \textbf{que} ampliem o letramento matemático iniciado na etapa anterior.
Segundo a BNCC, “ isso significa \textbf{que} novos conhecimentos específicos devem estimular processos mais elaborados de reflexão e de abstração ” ( BRASIL, 2018a, p. 529 ).
O alcance desse propósito ocorrerá quando os estudantes do Ensino Médio \textbf{conseguirem} atingir objetivos de aprendizagem \textbf{que} abordem processos de investigação, de construção de modelos e de resolução de problemas. \textbf{Contudo}, isso só será possível se os \textbf{discentes} forem estimulados a desenvolver o próprio modo de raciocinar, \textbf{que} está diretamente ligado a ser capaz de refletir e \textbf{argumentar}, individual e coletivamente ; aprendendo, assim, procedimentos cada vez mais sofisticados.
Desse modo, a metodologia para o ensino da Matemática deve ser centrada no estudante e, portanto, oferecer -lhe experiências pedagógicas nas quais ele possa refletir sobre seu próprio desempenho.
Por exemplo, verificar se a equação \textbf{que} elaborou em determinada atividade é adequada ou não, pois essa confirmação estará na própria Matemática.
Ao explicar o \textbf{que} fez, ele se organiza mentalmente, procura ordenar o seu ponto de vista e pode compartilhá -lo com o colega.
Ao discutir, ele é capaz de entender seu próprio \textbf{raciocínio} \textbf{enquanto} constrói seu conhecimento e relaciona -se com suas estruturas mentais e com o mundo físico e social.
De acordo com Mora ( 2003 ), > A Matemática somente será entendida, aprendida e dominada pela maioria das pessoas, quando sua relação com elas estiver \textbf{baseada}, em primeiro lugar, no trabalho ativo, participativo e significativo dos \textbf{sujeitos} \textbf{atores} do processo educativo ( MORA, 2003, p. 49 ).
Em consonância com o \textbf{autor}, D’Ambrosio alerta para a necessidade de \textbf{que} os professores compreendam a Matemática como \textbf{um} componente curricular em \textbf{que} o avanço da aprendizagem ocorrerá como consequência de \textbf{um} processo de investigação e resolução de problemas ( D’AMBROSIO, 1993 ).
Além disso, ressalta a importância da compreensão, por parte do docente, de \textbf{que} a Matemática a ser estudada deve ser útil aos estudantes, auxiliando -os a compreender e interferir em sua realidade.
Nessa perspectiva, as rápidas transformações na dinâmica social e \textbf{contemporânea}, nacional e internacional, em grande parte decorrentes do desenvolvimento tecnológico, atingem diretamente as populações jovens e, portanto, suas demandas de formação.
Eles costumam ser os mais \textbf{influenciados} pelas inovações tecnológicas, pois nascem e crescem interagindo com \textbf{um} mundo \textbf{que}, para \textbf{uma} grande quantidade de adultos, é \textbf{novidade}.
A educação, pela dinamicidade de sua natureza, transforma -se e \textbf{aprimora} -se, visando acompanhar essas mudanças e atender aos \textbf{imperativos} determinados pela vida \textbf{moderna}, na qual a tecnologia é usada diariamente.
Sendo assim, a tecnologia vinculada ao processo de ensino e aprendizagem torna -se \textbf{poderosa} ferramenta para a construção de conhecimentos no âmbito da Matemática.
Sobre isso, o PCNEM+ ressalta \textbf{que} \textbf{uma} competência a ser desenvolvida pelo estudante é a de “ perceber o papel \textbf{desempenhado} pelo conhecimento matemático no desenvolvimento da tecnologia e a complexa relação entre ciência e tecnologia ao longo da história ” ( BRASIL, 2002a, p. 118 ).
Desde os Anos Iniciais do Ensino Fundamental, já é proposto aos estudantes o uso de tecnologias, como calculadoras e planilhas eletrônicas, para a resolução de situações-problema.
Tal valorização possibilita \textbf{que} ao chegarem aos Anos Finais, eles possam ser estimulados a desenvolver o pensamento computacional, por meio da interpretação e da elaboração de algoritmos, incluindo aqueles \textbf{que} podem ser representados por fluxogramas.
Nesse contexto, para a etapa do Ensino Médio, enfatiza -se, nesta edição do Currículo em Movimento, a importância das tecnologias digitais para o ensino da Matemática - em especial, no \textbf{que} \textbf{diz} respeito à investigação matemática e ao desenvolvimento do pensamento computacional.
O uso da tecnologia integrada ao currículo pode ser \textbf{um} importante subsídio para o ensino, buscando sanar as dificuldades de aprendizagem, \textbf{uma} vez \textbf{que} possibilita torná -lo mais atrativo e significativo para os estudantes.
No entanto, o uso das tecnologias, por si só, não garante \textbf{um} ensino inovador, pois elas também podem reproduzir processos formais e repetitivos de aprendizagem.
É necessário \textbf{que} exista \textbf{uma} identificação cultural e, também, \textbf{que} o professor vislumbre a possibilidade de obter algum ganho no seu fazer pedagógico ( PONTE, 1992 ).
O uso de tecnologias vem para auxiliar a mediação pedagógica, aumentando, assim, a interatividade entre professor e estudante.
Dessa forma, o estudante se percebe como sujeito no seu processo de construção do conhecimento, e o professor assume o papel de mediador desse processo, promovendo a formação de valores, tais como : iniciativa, motivação, autodisciplina e autonomia.
A seguir, são apresentados os objetivos gerais da área de Matemática e suas Tecnologias, a organização curricular, as unidades temáticas e a estrutura utilizada para a escrita dos objetivos de aprendizagem.
A área de Matemática e suas Tecnologias visa, principalmente, possibilitar o desenvolvimento do \textbf{raciocínio} lógico-matemático e argumentativo do estudante, a fim de mostrar \textbf{que} o processo do descobrimento matemático é algo vivo e em desenvolvimento.
Para tanto, é necessário traçar \textbf{um} conjunto de objetivos \textbf{que} permitem colocar em prática essa premissa e subsidiar \textbf{um} planejamento interdisciplinar entre as áreas do conhecimento.
Os objetivos gerais elencados a seguir não estão organizados em \textbf{uma} ordem predeterminada.
Ou seja, todos possuem igual nível de importância e se conectam, de modo \textbf{que} o trabalho de desenvolvimento de \textbf{um} pressupõe o envolvimento dos demais.
Eles foram elaborados com base na interpretação das competências da área de Matemática e suas Tecnologias elencadas na BNCC.
Destaca -se \textbf{que} esses objetivos primam pelo desenvolvimento da autoestima do estudante, levando em consideração a busca individual pelo conhecimento e o respeito ao trabalho construído em grupo.
Sendo assim, ao final da etapa do Ensino Médio, na área de Matemática e suas Tecnologias, espera -se \textbf{que} os estudantes, de modo geral, sejam capazes de : - aplicar conceitos e procedimentos matemáticos a situações variadas, utilizando -os na interpretação de diversos contextos : atividades cotidianas, fatos das Ciências da Natureza e Humanas, questões socioeconômicas e tecnológicas, de modo a contribuir para \textbf{uma} formação integral ; - mobilizar e articular conceitos, procedimentos e linguagens próprios da Matemática, a fim de propor ações e soluções para problemas sociais, com base na investigação do mundo \textbf{contemporâneo}, tomando decisões éticas e socialmente responsáveis ; - resolver e elaborar situações-problema em diversos contextos, incluindo os oriundos do desenvolvimento tecnológico, analisando a plausibilidade dos resultados e a adequação das soluções propostas, de modo a construir \textbf{uma} argumentação consistente ; - utilizar diferentes registros matemáticos de representação ( algébrico, geométrico, estatístico, computacional, entre outros ) na resolução de problemas em diferentes contextos, como por exemplo, os socioambientais e os da vida cotidiana, escolhendo as representações mais convenientes a cada situação e convertendo -as sempre \textbf{que} necessário ; - elaborar conjecturas a respeito de diferentes conceitos e propriedades matemáticas, identificando a necessidade, ou não, de \textbf{uma} demonstração cada vez mais formal na validação, utilizando estratégias e recursos, como observação de padrões, experimentações e diferentes tecnologias.

% matched lemmas: aprimorar, argumentar, ator, atual, autor, basear, conseguir, contemporâneo, contudo, desempenhar, discente, dizer, enquanto, imperativo, influenciar, lembrar, memorização, moderno, novidade, permear, poderoso, popular, predominante, qualificar, que, raciocínio, simples, sujeito, século, um, ótica
\end{textsample}
